% !PRINTED 511
\weberpage{511}

We therefore call a divisor $Q$ of the group $P$ which has this
property a \gesperrt{normal divisor} (or normal factor)\footnote{\gesperrt{Galois} speaks of the proper decomposition (d\'ecomposition propre) of a group; hence the expression ``proper divisor'' which is also in use; more recent writers designate the normal divisors as ``distinguished or invariant subgroups''.}.

But if $Q$ is a normal divisor, then the adjunction of a
single root $\psi$ of the auxiliary equation is equivalent to the
simultaneous adjunction of all these roots.

5. \gesperrt{If $Q$ itself is not normal, then the greatest
common divisor $R$ of all the conjugate
divisors $\pi^{-1}Q\pi$ is certainly normal.}

For if a permutation $\varkappa$ is contained in $R$, and hence in all the groups
$\pi^{-1}Q\pi$, then the same holds also of every $\pi^{-1}\varkappa\pi$;
that is, $\pi^{-1}R\pi$ is a divisor of $R$, and consequently, since it is of the same
degree, $= R$.

A group which has no normal divisor other than itself and the identical
group is called a \gesperrt{simple group}.
The intersection of all the groups conjugate to a simple group
is either the group itself or the identical group.

If $Q$ is a normal divisor of $P$, and $R$ a normal divisor
of $Q$, then $R$ is indeed a divisor of $P$, but by no means always
a normal one. On the other hand, if $R$ is a normal divisor of $P$, then $R$
is also a normal divisor of every divisor $Q$ of $P$ of which $R$
is a divisor.

\begin{center}
\S.\,156.

\gesperrt{The Group of the Resolvents.}
\end{center}

The auxiliary equation $\varphi(t) = 0$, on which the determination of the
function $\psi$ depends, passes over into the \gesperrt{Galois} resolvent
when the divisor $Q$ to which $\psi$ belongs is the identical group.
For then, by the theorem of \gesperrt{Lagrange} (\S.\,155, 2.),
every function of the field $N$, and hence also the roots $\alpha$ themselves,
can be expressed rationally in terms of $\psi$, and $N$ is identical with $\Omega(\psi)$.

We shall therefore call these equations $\varphi(t) = 0$, in a general
sense, \gesperrt{resolvents}. But two cases are here
to be distinguished.
